% exam1-prop-logic.tex

\newpage
\section{测验 1：命题逻辑}

\begin{problem}[命题逻辑的自然推理系统]
  请使用自然推理规则证明以下结论:

  \begin{enumerate}[(1)]
    \setcounter{enumi}{3}
    \item
      \[
        \vdash ((p \to q) \to q) \to ((q \to p) \to p)
      \]
      \begin{solution}
        \begin{intention}
          外层连做两次 $\to\intro$; 真正的难点是从 $\lnot p$ 构造出 $p \to q$。
        \end{intention}

        {\small
        \begin{logicproof}{4}
          \begin{subproof}
            (p \to q) \to q & assumption \\
            \begin{subproof}
              q \to p & assumption \\
              \begin{subproof}
                \lnot p & assumption \\
                \begin{subproof}
                  p & assumption \\
                  \bot & $\lnot\elim$ 3, 4 \\
                  q & $\bot\elim$ 5
                \end{subproof}
                p \to q & $\to\intro$ 4--6 \\
                q & $\to\elim$ 1, 7 \\
                p & $\to\elim$ 2, 8 \\
                \bot & $\lnot\elim$ 3, 9
              \end{subproof}
              \lnot\lnot p & $\lnot\intro$ 3--10 \\
              p & $\lnot\lnot\elim$ 11
            \end{subproof}
            (q \to p) \to p & $\to\intro$ 2--12
          \end{subproof}
          ((p \to q) \to q) \to ((q \to p) \to p) & $\to\intro$ 1--13
        \end{logicproof}}

        \begin{remark}
          \textbf{难点.} 在假设 $\lnot p$ 下, 通过子证明
          $p \Rightarrow \bot \Rightarrow q$
          得到 $p \to q$, 这是整题的关键构造。
          \textbf{易错点.} 两层 $\to\intro$ 与一层 $\lnot\intro$ 的释放区间不要写错。
        \end{remark}
      \end{solution}
  \end{enumerate}
\end{problem}

%%%%%%%%%%%%%%%

\begin{problem}[命题逻辑的自然推理系统]
  请使用自然推理规则证明以下结论:

  \begin{enumerate}[(1)]
    \setcounter{enumi}{2}
    \item (德摩根律)
      \[
        \lnot(\alpha \land \beta) \vdash (\lnot\alpha \lor \lnot\beta)
      \]
      提示: 你可能需要多次用到排中律 $\frac{}{\alpha \lor \lnot\alpha}$。
      \begin{solution}
        \begin{intention}
          这是\blue{经典}结论; 按提示两次使用排中律, 再做分情况讨论。
        \end{intention}

        {\small
        \begin{logicproof}{4}
          \lnot(\alpha \land \beta) & premise \\
          \alpha \lor \lnot\alpha & LEM \\
          \begin{subproof}
            \alpha & assumption \\
            \beta \lor \lnot\beta & LEM \\
            \begin{subproof}
              \beta & assumption \\
              \alpha \land \beta & $\land\intro$ 3, 5 \\
              \bot & $\lnot\elim$ 1, 6 \\
              \lnot\alpha \lor \lnot\beta & $\bot\elim$ 7
            \end{subproof}
            \begin{subproof}
              \lnot\beta & assumption \\
              \lnot\alpha \lor \lnot\beta & $\lor\introright$ 9
            \end{subproof}
            \lnot\alpha \lor \lnot\beta & $\lor\elim$ 4, 5--8, 9--10
          \end{subproof}
          \begin{subproof}
            \lnot\alpha & assumption \\
            \lnot\alpha \lor \lnot\beta & $\lor\introleft$ 12
          \end{subproof}
          \lnot\alpha \lor \lnot\beta & $\lor\elim$ 2, 3--11, 12--13
        \end{logicproof}}

        \begin{remark}
          \textbf{难点.} 该结论不是纯直觉主义写法; 这里明确依赖排中律 LEM。
          \textbf{易错点.} 在情形 $\alpha, \beta$ 下, 先得出 $\alpha \land \beta$, 再与前提碰撞得到 $\bot$,
          最后用 $\bot\elim$ 返回目标式。
        \end{remark}
      \end{solution}
  \end{enumerate}
\end{problem}

%%%%%%%%%%%%%%%

\begin{problem}[另一个命题逻辑推理系统]
  考虑一个``希尔伯特''式的命题逻辑推理系统,
  它包含三个公理 (A1)、(A2) 与 (A3):
  \begin{enumerate}[label=(A\arabic*)]
    \item
      \[
        \alpha \to (\beta \to \alpha)
      \]
    \item
      \[
        (\alpha \to (\beta \to \gamma)) \to ((\alpha \to \beta) \to (\alpha \to \gamma))
      \]
    \item
      \[
        (\lnot \beta \to \lnot \alpha) \to ((\lnot \beta \to \alpha) \to \beta)
      \]
  \end{enumerate}
  除此之外, 它还包含一个推理规则 (与课堂上介绍的自然推理系统不同,
  它只包含一个推理规则):
  \[
    \frac{\qquad\alpha \to \beta \qquad \alpha\qquad}{\beta} \qquad \text{MP; 分离规则}
  \]

  在该推理系统中, 公式 $\alpha \to \alpha$ 的一种证明方式如下:
  \begin{align}
    (\alpha \to (\alpha \to \alpha) \to \alpha) \to
      (\alpha \to (\alpha \to \alpha)) \to (\alpha \to \alpha) \\
    \alpha \to (\alpha \to \alpha) \to \alpha \\
    (\alpha \to (\alpha \to \alpha)) \to (\alpha \to \alpha) \\
    \alpha \to (\alpha \to \alpha) \\
    \alpha \to \alpha
  \end{align}
  请给出每一步的依据 (即, 基于哪几条公式, 应用哪条公理或分离规则)。
\end{problem}

\begin{solution}
  \begin{intention}
    先识别公理模式的代换实例, 再连续两次应用 MP。
  \end{intention}

  \begin{align}
    &(\alpha \to (\alpha \to \alpha) \to \alpha) \to
      (\alpha \to (\alpha \to \alpha)) \to (\alpha \to \alpha)
      && \text{A2, 令 $\beta := (\alpha \to \alpha),\ \gamma := \alpha$} \tag{1} \\
    &\alpha \to (\alpha \to \alpha) \to \alpha
      && \text{A1, 令 $\beta := (\alpha \to \alpha)$} \tag{2} \\
    &(\alpha \to (\alpha \to \alpha)) \to (\alpha \to \alpha)
      && \text{MP, 由 (1), (2)} \tag{3} \\
    &\alpha \to (\alpha \to \alpha)
      && \text{A1, 令 $\beta := \alpha$} \tag{4} \\
    &\alpha \to \alpha
      && \text{MP, 由 (3), (4)} \tag{5}
  \end{align}

  \begin{remark}
    \textbf{重点.} 第 (1) 步是 A2 的代换实例, 不是 A1。
    \textbf{易错点.} A1 的固定形状是 $\alpha \to (\beta \to \alpha)$; 代换时要先确定“首尾同一个 $\alpha$”。
  \end{remark}
\end{solution}
